\chapter{Metodologia}
\label{cap:metodologia}

\section{O Modelo FactorVAE}
\label{sec:modelo-factorvae}

\subsection{Formulação e Notação}
\label{subsec:formulacao-notacao}

\begin{itemize}
    \item \textbf{Cross-section e notação:} em cada data $t$, há $N_t$ ativos disponíveis ($\approx 300$--$500$ no período de teste); o vetor de retornos realizados é $r_t \in \mathbb{R}^{N_t}$; o tensor de características históricas é $x_t \in \mathbb{R}^{N_t \times T \times C}$, com $T = 20$ dias de janela temporal e $C = 20$ características por ativo por período.
    \item \textbf{Variável-alvo:} retorno logarítmico de $t{+}1$ para $t{+}2$ --- um passo à frente relativo ao fechamento observado em $t$; o alinhamento temporal elimina look-ahead (detalhes na Subseção~\ref{subsec:universo-periodo}).
    \item \textbf{Equação geral do modelo dinâmico de fatores:}
\end{itemize}

\begin{equation}
    r_{t+1} = \alpha(x_t) + \beta(x_t)\, z_{t+1} + \varepsilon_{t+1}
    \label{eq:decomp-fatorial}
\end{equation}

\begin{itemize}
    \item[] em que $\alpha(x_t) \in \mathbb{R}^{N_t}$ é o vetor de componentes idiossincráticos esperados, $\beta(x_t) \in \mathbb{R}^{N_t \times K}$ é a matriz de exposições fatoriais, $z_{t+1} \in \mathbb{R}^K$ é o vetor de fatores latentes e $\varepsilon_{t+1}$ é o resíduo idiossincrático. Todas as três quantidades são funções aprendidas de $x_t$.
    \item \textbf{Parâmetros de dimensão:} $K = 16$ fatores latentes, $M = 64$ portfólios internos do codificador, dimensão oculta $H$ especificada no Apêndice~\ref{ap:hiperparametros}. Escolhas seguem \textcite{duan2022} para comparabilidade direta.
    \item \textbf{Pesos compartilhados entre ativos:} todas as redes processam cada ativo de forma independente com pesos compartilhados --- isso é o que permite generalizar para ativos não vistos durante o treinamento, conforme avaliado na Subseção~\ref{subsec:avaliacao-robustez}.
\end{itemize}

\subsection{Caminho de Inferência}
\label{subsec:caminho-inferencia}

O caminho de inferência é o fluxo que opera fora da amostra, com acesso apenas a informação disponível até $t$.

\begin{itemize}
    \item \textbf{Histórico $\to$ embeddings (Extrator de Características):} para cada ativo $i$, uma projeção linear seguida de GRU \parencite{cho2014} processa a janela $x_t^{(i)} \in \mathbb{R}^{T \times C}$ e produz o embedding $e^{(i)} \in \mathbb{R}^H$ --- representação compacta do padrão temporal relevante; pesos compartilhados entre ativos.
    \item \textbf{Embeddings $\to$ preditor (Preditor de Fatores):} $K = 16$ cabeças de atenção independentes agregam os embeddings com pesos baseados em similaridade cosseno (gate ReLU) e projetam nos parâmetros da distribuição prior $p(z|x_t) = \mathcal{N}(\mu_{\text{prior}}, \mathrm{diag}(\sigma_{\text{prior}}^2))$.
    \item \textbf{Preditor $\to$ decodificador $\to$ distribuição preditiva (Decodificador):} os parâmetros $\mu_{\text{prior}}$ e $\sigma_{\text{prior}}$ são combinados com os embeddings para produzir a distribuição preditiva por ativo. Sob independência condicional gaussiana, média e variância têm forma fechada:
\end{itemize}

\begin{align}
    \mu_y^{(i)} &= \mu_\alpha^{(i)} + \beta^{(i)\top}\, \mu_{\text{prior}} \label{eq:mu-preditivo} \\
    \sigma_y^{(i)} &= \sqrt{\left(\sigma_\alpha^{(i)}\right)^2 + \sum_{k=1}^{K} \left(\beta^{(i,k)}\right)^2 \left(\sigma_{\text{prior}}^{(k)}\right)^2} \label{eq:sigma-preditivo}
\end{align}

\begin{itemize}
    \item \textbf{Estimativa de risco como subproduto:} $\sigma_y^{(i)}$ (Equação~\ref{eq:sigma-preditivo}) é a incerteza preditiva por ativo entregue em forma fechada, sem Monte Carlo --- aproveitável na construção de estratégias ajustadas ao risco (Capítulo~\ref{cap:conclusao}).
    \item \textbf{O codificador não participa da inferência:} o codificador depende dos retornos contemporâneos $y_t$, indisponíveis no momento da previsão; existe apenas durante o treinamento.
\end{itemize}

% PLACEHOLDER: figura do fluxo de inferência --- diagrama mostrando os três blocos ativos
% na inferência (Extrator -> Preditor -> Decodificador) com indicação das dimensões dos
% tensores em cada etapa. Elaboração própria.

\subsection{Caminho de Treino}
\label{subsec:caminho-treino}

\begin{itemize}
    \item \textbf{Codificador (com acesso a retornos contemporâneos):} recebe os embeddings $\{e^{(i)}\}$ e os retornos contemporâneos $y_t \in \mathbb{R}^{N_t}$; uma \textit{PortfolioLayer} constrói $M = 64$ portfólios cujos pesos dependem dos embeddings (não dos retornos); a \textit{MappingLayer} projeta os retornos dos portfólios nos parâmetros da posterior $q(z|x_t, y_t) = \mathcal{N}(\mu_{\text{post}}, \mathrm{diag}(\sigma_{\text{post}}^2))$.
    \item \textbf{Esquema prior-posterior:} durante o treino, a posterior (codificador, informada pelos retornos contemporâneos) funciona como professor; a prior (preditor, apenas histórico) funciona como aluno. A função objetivo força a prior a se aproximar da posterior --- mecanismo que transfere o sinal preditivo do oráculo ao preditor.
    \item \textbf{Função objetivo (negativo do ELBO adaptado):}
\end{itemize}

\begin{equation}
    \mathcal{L}(x, y) = \underbrace{-\frac{1}{N}\sum_{i=1}^{N} \log p_{\phi_{\text{dec}}}\!\left(y^{(i)} \;\middle|\; x,\, z_{\text{post}}\right)}_{\text{NLL (qualidade de reconstrução)}} + \;\gamma \cdot \underbrace{\mathrm{KL}\!\left(q_{\phi_{\text{enc}}}(z|x,y) \;\|\; p_{\phi_{\text{pred}}}(z|x)\right)}_{\text{regularização (prior-posterior)}}
    \label{eq:loss}
\end{equation}

\begin{itemize}
    \item \textbf{Interpretação dos termos:} NLL (média sobre $N$ ativos) mede a qualidade da previsão dos retornos contemporâneos via posterior; KL (média sobre $K$ dimensões latentes) força o preditor a aprender fatores consistentes com o oráculo. A função objetivo é consequência da arquitetura, não ponto de partida.
    \item \textbf{Hiperparâmetro $\gamma$:} controla o peso relativo da regularização; agendamento linear durante os primeiros 500 passos (\textit{KL warmup}) e decisões de estabilização numérica documentados no Apêndice~\ref{ap:estabilizacao}.
    \item \textbf{Otimizador:} Adam, \textit{learning rate} $3 \times 10^{-4}$, \textit{weight decay} $10^{-4}$; \textit{batch} de uma cross-section por passo; \textit{early stopping} monitorando Rank IC de validação com paciência de 10 épocas.
    \item \textbf{Validade do esquema:} o diagnóstico de Posterior IC (Apêndice~\ref{ap:diagnosticos}) verifica que o codificador captura sinal preditivo genuíno dos retornos contemporâneos --- não apenas reconstrói a entrada trivialmente.
\end{itemize}

\section{Dados}
\label{sec:dados}

\subsection{Universo, Período e Variável-Alvo}
\label{subsec:universo-periodo}

\begin{itemize}
    \item \textbf{Critérios de inclusão na B3:} todos os tickers com \textit{pricing} válido na Economatica, sem restrição a índices de referência (IBX ou similares); ticker incluído na cross-section da data $t$ se e somente se apresenta $T = 20$ dias úteis consecutivos sem lacuna calendárica superior a 7 dias (tolerância para feriados nacionais); preço mínimo de R\$~0{,}10 (exclui penny stocks); retorno diário absoluto máximo de 50\% (exclui eventos de split não ajustado e erros de dados).
    \item \textbf{Tamanho do universo:} $\approx 300$--$500$ ativos no período de teste --- substancialmente menor do que o mercado A-share chinês original ($>3.000$); diferença estrutural que condiciona a interpretação de todas as métricas de portfólio e a comparação de magnitudes com o artigo original.
    \item \textbf{Partição treino / validação / teste (sem sobreposição temporal):}
    % PLACEHOLDER: definir datas exatas após consolidar o dataset.
    % Estrutura esperada:
    %   Treinamento: [data_inicio] a [data_fim_treino]
    %   Validação:   [data_inicio_val] a [data_fim_val]
    %   Teste:       [data_inicio_teste] a [data_fim_teste]
    % O período de teste deve cobrir pelo menos: pandemia 2020, aperto monetário 2021--2022,
    % mudança de governo 2023.
    \item \textbf{Variável-alvo e cuidados contra look-ahead:} retorno logarítmico $r_{t+1} = \ln(P_{t+2}/P_{t+1})$ --- calculado com preços de fechamento dois dias à frente; o modelo treinado na data $t$ recebe como entrada apenas características calculadas até $t$; as estatísticas de normalização (médias e desvios-padrão) são calculadas exclusivamente sobre o período de treinamento e aplicadas sem re-estimação aos períodos de validação e teste.
    \item \textbf{Auditoria de viés de sobrevivência:} a construção point-in-time reconstrói a composição do universo disponível em cada data histórica sem informação prospectiva, preservando tickers que deslistaram durante o período amostral; protocolo documentado no Apêndice~\ref{ap:diagnosticos}.
\end{itemize}

\subsection{Características e Pré-processamento}
\label{subsec:caracteristicas-preprocessing}

\begin{itemize}
    \item \textbf{Conjunto de $C = 20$ características técnicas:} calculadas a partir de preço e volume, organizadas em seis categorias --- retornos (5), volatilidade (4), volume (3), indicadores técnicos (3), momentos superiores (2) e iliquidez (1); justificativa para o subconjunto: disponibilidade na Economatica, representatividade das categorias de sinal e alinhamento com \textcite{duan2022}. Fórmulas detalhadas no Apêndice~\ref{ap:caracteristicas}.
\end{itemize}

\begin{table}[htbp]
    \centering
    \caption{Características utilizadas como entrada do modelo.}
    \label{tab:caracteristicas}
    \begin{tabular}{llc}
        \toprule
        \textbf{Nome} & \textbf{Categoria} & \textbf{Janela (dias)} \\
        \midrule
        Ret1, Ret2, Ret5, Ret10, Ret20  & Retorno         & 1, 2, 5, 10, 20 \\
        Vol5, Vol10, Vol20              & Volatilidade    & 5, 10, 20       \\
        VolRatio                        & Volatilidade    & 5/20            \\
        Turnover5, Turnover10           & Volume          & 5, 10           \\
        VWAPDev                         & Volume          & 5               \\
        RSI                             & Indicador téc.  & 14              \\
        MADev5, MADev20                 & Indicador téc.  & 5, 20           \\
        Skew20, Kurt20                  & Momentos sup.   & 20              \\
        Amihud20                        & Iliquidez       & 20              \\
        \bottomrule
    \end{tabular}
    \fonte{Elaboração própria. Fórmulas detalhadas no Apêndice~\ref{ap:caracteristicas}.}
\end{table}

\begin{itemize}
    \item \textbf{Winsorização:} cada característica é winsorizada cross-sectionalmente a $[p_1, p_{99}]$ a cada data --- valores extremos substituídos pelos percentis 1 e 99 do cross-section naquela data; evita que outliers distorçam os embeddings.
    \item \textbf{Normalização cross-sectional:} \textit{z-score} sobre os $N_t$ ativos disponíveis na data $t$, calculado separadamente para cada característica; as estatísticas de normalização são estimadas apenas sobre o período de treinamento e aplicadas sem re-estimação a validação e teste.
\end{itemize}

\section{Procedimento de Avaliação}
\label{sec:avaliacao}

\subsection{Modelos Comparados}
\label{subsec:modelos-comparados}

A progressão de baselines isola sucessivamente as propriedades da arquitetura do \textit{FactorVAE}.

\begin{itemize}
    \item \textbf{EW Market (benchmark de referência):} portfólio igual-ponderado sobre o universo de ativos disponíveis em cada data --- sem sinal preditivo; serve como benchmark de portfólio para calcular excesso de retorno e Information Ratio.
    \item \textbf{Momentum:} retorno dos últimos 20 dias (Ret20) como previsão direta --- o sinal econômico mais simples, sem aprendizado; representativo do prêmio de momento documentado na literatura.
    \item \textbf{Ridge:} regressão Ridge com as mesmas $C = 20$ características como entrada --- limite do problema sob hipótese de linearidade e regularização L2; avalia se a não-linearidade das redes neurais agrega valor.
    \item \textbf{MLP:} rede \textit{feedforward} com as mesmas características --- acrescenta não-linearidade sobre o Ridge, sem estrutura temporal; avalia o ganho de representações não-lineares.
    \item \textbf{GRU:} rede recorrente com as mesmas características --- acrescenta estrutura temporal sobre o MLP, sem componente probabilístico e sem estrutura fatorial explícita; representa a alternativa neural mais próxima do FactorVAE em complexidade.
    \item \textbf{Nota metodológica:} os modelos diferem do \textit{FactorVAE} em mais de uma dimensão simultaneamente; a comparação não constitui ablação formal --- a atribuição causal estrita ficaria a cargo de um estudo de ablação dedicado, fora do escopo deste trabalho.
\end{itemize}

\subsection{Métricas e Construção de Carteira}
\label{subsec:avaliacao-robustez}

\begin{itemize}
    \item \textbf{Rank IC:} correlação de Spearman entre retornos previstos e retornos realizados no \textit{cross-section} de cada data $t$; média ao longo do período de teste. Mede a qualidade do sinal preditivo sem depender da estratégia de portfólio.
    \item \textbf{Rank ICIR:} razão entre a média e o desvio-padrão das observações diárias de Rank IC --- mede a consistência temporal do sinal; sinal com alto IC mas ICIR baixo é instável e difícil de monetizar.
    \item \textbf{Estratégia TopK-Drop \parencite{yang2020}:} a cada data, selecionar os $k = 50$ ativos com maior retorno previsto; rebalancear a carteira com no máximo $n = 5$ substituições por dia (controla turnover); portfólio igual-ponderado entre os $k$ ativos selecionados.
    \item \textbf{Adaptação ao universo brasileiro:} $k = 50$ representa $\approx 15$--$25\%$ do \textit{cross-section} típico no período de teste --- fração substancialmente maior do que os $\approx 1{,}5\%$ do mercado chinês; a comparação direta de magnitudes de retorno é problemática e deve ser interpretada com cautela.
    \item \textbf{Custo de transação:} 10~bps \textit{one-way} por posição substituída ($\approx 20$~bps \textit{round-trip}), aplicado a cada troca de ativo a cada dia; custo constante --- não captura impacto de mercado em ativos menos líquidos (limitação discutida no Capítulo~\ref{cap:conclusao}).
    \item \textbf{Métricas de portfólio:} retorno anualizado, retorno em excesso sobre o benchmark EW Market, volatilidade anualizada, Índice de Sharpe, Information Ratio, Índice de Calmar, drawdown máximo e hit rate (fração de períodos com retorno positivo).
    \item \textbf{Robustez a ativos não vistos (holdout-retrain):} $m$ tickers selecionados aleatoriamente são removidos do treinamento; o modelo é retreinado do zero; Rank IC é calculado apenas sobre esses $m$ tickers no período de teste. Procedimento repetido para $m \in \{10, 50\}$ com múltiplas amostras aleatórias --- quantifica a capacidade de generalização a IPOs e ativos de baixa liquidez histórica.
\end{itemize}
